\documentclass[11pt]{article}
\usepackage[margin=1in]{geometry}
\usepackage{amsmath,amssymb,amsthm,mathtools}
\usepackage{graphicx}
\usepackage{float}
\usepackage{microtype}
\usepackage[hidelinks]{hyperref}

\newtheorem{theorem}{Theorem}
\newtheorem{lemma}[theorem]{Lemma}
\newcommand{\R}{\mathbb R}
\newcommand{\gauss}{\gamma}

\title{Gaussian Balls Refute the Proposed\\Complement-Concavity Extension}
\author{Counterexample packet for arXiv:2407.20064}
\date{12 August 2026}

\begin{document}
\maketitle

\begin{center}
\fbox{\begin{minipage}{0.9\textwidth}
\textbf{Status: candidate full counterexample, likely valid; human review
requested.}
The source's final conjecture is false already for the standard Gaussian in
dimension two.  With \(K=6B_2^2\), \(L=2B_2^2\), and \(t=1/2\), its left
side is \(e^{-2}\), while its right side is strictly larger than
\((1-e^{-2})/4>e^{-2}\).  More generally, every absolutely continuous
probability measure satisfying the conjecture's concavity premise violates
its proposed conclusion for some two concentric balls.
\end{minipage}}
\end{center}

\section{The conjecture}

Let \(\mathcal K_e^n\) denote the origin-symmetric convex bodies in
\(\R^n\).  The final conjecture in Langharst--Liu--Tang,
\emph{The Weighted \(L^p\) Minkowski Problem} \cite{LLT}, says that if a
probability measure \(\mu\) is \((1/n)\)-concave over
\(\mathcal K_e^n\), then for all \(K,L\in\mathcal K_e^n\) and
\(t\in(0,1)\),
\[
\begin{split}
 &\mu\bigl((1-t)(\R^n\setminus\operatorname{int}K)+tL\bigr)\\
 &\quad\ge
 \left((1-t)\mu(\R^n\setminus\operatorname{int}K)^{1/n}
       +t\mu(L)^{1/n}\right)^n.
\end{split}
\tag{1}
\]

\begin{figure}[H]
\centering
\includegraphics[width=0.97\textwidth]{figures/source_conjecture_crop.png}
\caption{The source's final conjecture (PDF page 51).}
\end{figure}

The premise means that the same inequality, without complements and with the
usual Minkowski combination, is already known for pairs of symmetric convex
bodies.  The proposed extension replaces one convex input by its nonconvex,
unbounded complement.  That change creates an elementary obstruction.

\section{A general ball obstruction}

Write \(B=B_2^n\) for the Euclidean unit ball.

\begin{lemma}[Exact complement sum]
\label{lem:set}
Let \(a>b>0\) and set \(t=(a-b)/(a+b)\).  Then
\[
 (1-t)(\R^n\setminus\operatorname{int}(aB))+tbB
 =\R^n\setminus\operatorname{int}(bB).
 \tag{2}
\]
\end{lemma}

\begin{proof}
The choice of \(t\) gives
\((1-t)a-tb=b\).  If \(|y|\ge(1-t)a\) and \(|z|\le tb\), the reverse
triangle inequality gives \(|y+z|\ge b\), proving one inclusion.

Conversely, take \(|x|\ge b\).  If \(|x|\ge(1-t)a\), use \(y=x,z=0\).
If \(b\le|x|<(1-t)a\), use
\(y=(1-t)a\,x/|x|\) and \(z=x-y\).  Then
\(|z|=(1-t)a-|x|\le(1-t)a-b=tb\), so \(x=y+z\) belongs to the left side.
\end{proof}

\begin{theorem}[Structural failure]
\label{thm:general}
Let \(\mu\) be an absolutely continuous probability measure on \(\R^n\)
which is \((1/n)\)-concave over \(\mathcal K_e^n\).  Then the conclusion
(1) fails for some \(K,L\in\mathcal K_e^n\) and \(t\in(0,1)\).
\end{theorem}

\paragraph{Proof intuition.}
Choose a ball \(L=bB\) carrying more than half the probability.  Take a much
larger concentric ball \(K=aB\) and tune \(t\) so that adding \(tL\) to the
scaled exterior of \(K\) produces exactly the exterior of \(L\).  The left
side of (1) is then the small mass outside \(L\), while the right side retains
almost all the large mass inside \(L\).  As \(a\to\infty\), the coefficient
of that impossible term tends to one.

\begin{proof}
Choose \(b>0\) with \(\alpha:=\mu(bB)>1/2\); absolute continuity makes
\(\mu(\partial(bB))=0\).  For \(a>b\), set
\(K=aB\), \(L=bB\), and \(t=(a-b)/(a+b)\).  Lemma~\ref{lem:set} makes the
left side of (1) equal to \(1-\alpha\).  Its right side is at least
\[
 \bigl(t\mu(L)^{1/n}\bigr)^n=t^n\alpha. \tag{3}
\]
As \(a\to\infty\), \(t\to1\).  Since \(\alpha>1/2\), one can choose a
finite \(a\) for which \(t^n\alpha>1-\alpha\).  For that choice, (3) is
strictly larger than the left side, contradicting (1).
\end{proof}

Notice that the concavity premise is used to make \(\mu\) an admissible
measure for the conjecture; the failure mechanism itself needs only
absolute continuity and finite total mass.

\section{An explicit Gaussian counterexample}

Let \(\gauss_2\) be standard Gaussian measure on \(\R^2\):
\[
 d\gauss_2(x)=\frac1{2\pi}e^{-|x|^2/2}\,dx.
\]
Eskenazis and Moschidis proved the dimensional Gaussian
Brunn--Minkowski inequality \cite{EM}: for symmetric convex bodies
\(A,C\subset\R^n\),
\[
 \gauss_n((1-s)A+sC)^{1/n}
 \ge(1-s)\gauss_n(A)^{1/n}+s\gauss_n(C)^{1/n}.
 \tag{4}
\]
Thus \(\gauss_2\) satisfies exactly the premise of (1).

Take
\[
 n=2,\qquad K=6B_2^2,\qquad L=2B_2^2,\qquad t=\frac12.
 \tag{5}
\]
Lemma~\ref{lem:set}, or a direct radius calculation \(3-1=2\), gives
\[
 \tfrac12(\R^2\setminus\operatorname{int}K)+\tfrac12L
 =\R^2\setminus\operatorname{int}(2B_2^2). \tag{6}
\]
Polar integration yields
\[
 \gauss_2(rB_2^2)=1-e^{-r^2/2}. \tag{7}
\]
Consequently, the left side of (1) is \(e^{-2}\), whereas the right side is
\[
 \frac14\left(e^{-9}+\sqrt{1-e^{-2}}\right)^2
 >\frac14(1-e^{-2}). \tag{8}
\]
Finally,
\[
 e>1+1+\frac12+\frac16=\frac83
 \quad\Longrightarrow\quad e^2>\frac{64}{9}>5,
\]
and hence
\[
 e^{-2}<\frac14(1-e^{-2}). \tag{9}
\]
Equations (8)--(9) show that the proposed right side is strictly larger than
the left side.  All bodies are origin-symmetric convex bodies, and (4)
verifies the only non-elementary hypothesis.

\section{Verification, scope, and novelty}

The counterexample is a full negative answer to the displayed source
conjecture, not to the weighted \(L^p\) Minkowski existence theorems proved
elsewhere in the paper.  It also does not challenge Gaussian
\((1/n)\)-concavity on symmetric convex bodies: one input in (1) is the
complement of such a body and is neither convex nor bounded.

Review should focus on the exact set identity (2), the Gaussian tail formula
(7), and the direction of the strict comparison (9).  The accompanying
script checks the radius, exact elementary certificate, and numerical gap.

On 12 August 2026, the four run indexes and bounded exact/current searches
for the arXiv id, title, exact complement formula, and later
citations found no existing proof or refutation.  The current source version
still states the conjecture on page 51.  Novelty confidence is moderate-high.

\paragraph{Human-review recommendation.}
Review as a full counterexample.  The example is explicit, two-dimensional,
uses the canonical measure motivating the paper, and extends to every
absolutely continuous probability measure satisfying the stated premise.

\begin{thebibliography}{9}

\bibitem{LLT}
D.~Langharst, J.~Liu, and S.~Tang,
\emph{The Weighted \(L^p\) Minkowski Problem},
arXiv:2407.20064, Appendix conjecture on PDF page 51.

\bibitem{EM}
A.~Eskenazis and G.~Moschidis,
\emph{The dimensional Brunn--Minkowski inequality in Gauss space},
J. Funct. Anal. 280 (2021), Paper No.~108914; arXiv:2004.07146.

\end{thebibliography}

\end{document}
